\documentclass{article}

\usepackage{amsmath}
\usepackage{listings}
\usepackage{amssymb}


\begin{document}


    \section{Theoretische Grundlagen}
    \label{sec:grundlagen}

    Dieses Kapitel legt die mathematischen und physikalischen Grundlagen dar,
    auf denen die Simulationspipeline aufbaut.
    Das Verständnis dieser Konzepte ist Voraussetzung für die Bewertung der
    in späteren Kapiteln beschriebenen Implementierungsentscheidungen.

    \subsection{Akustische Grundlagen}
    \label{subsec:akustik}

    \paragraph{Das Inverse Quadratgesetz}

    Im freien, reflexionsfreien Schallfeld breitet sich Schall kugelförmig vom Emitter aus.
    Da die abgestrahlte Schallleistung $P$ über eine mit $r^2$ wachsende Kugeloberfläche
    $A = 4\pi r^2$ verteilt wird, nimmt die Schallintensität $I$ mit dem Quadrat der Entfernung ab:
    \[I(r) = \frac{P}{4\pi r^2}\]
    Da die Schallintensität proportional zum Quadrat des Schalldrucks ist ($I \propto p^2$),
    folgt für den Schalldruck:
    \[p(r) \propto \frac{1}{r}\]
    Eine Verdopplung des Abstands zur Quelle halbiert den Schalldruck, was einer
    Pegelminderung von $6\,\text{dB}$ entspricht.
    In der Simulation wird dieser geometrische Energieverlust implizit berücksichtigt:
    Strahlen, die einen längeren Weg zur Mikrofonkapsel zurücklegen, treffen mit einem
    entsprechend geringeren Druckbeitrag ein.

    \paragraph{Materialabsorption}

    Trifft ein Schallstrahl auf eine Raumgrenze, wird ein Teil der Schallenergie absorbiert.
    Dieser Vorgang wird durch den dimensionslosen Absorptionskoeffizienten $\alpha \in [0{,}1]$
    beschrieben. Der Restdruck nach einer einzelnen Reflexion berechnet sich zu:
    \[p_{\text{ref}} = p_{\text{ein}} \cdot (1 - \alpha)\]
    Bei einer Sequenz von $n$ Reflexionen mit den jeweiligen Koeffizienten
    $\alpha_1, \ldots, \alpha_n$ gilt:
    \[p_n = p_0 \cdot \prod_{k=1}^{n}(1 - \alpha_k)\]
    Ein vollständig absorbierendes Material ($\alpha = 1{,}0$) vernichtet den Strahl bei der
    ersten Reflexion, während ein perfekt reflektierendes Material ($\alpha = 0{,}0$) keine
    Energie entzieht. Die Simulation terminiert Strahlen, sobald ihr verbleibender Druckwert
    den konfigurierbaren Grenzwert \texttt{min\_pressure} unterschreitet.

    \paragraph{Frühe Reflexionen und späte Nachhallzeit}

    Die akustische Impulsantwort eines Raumes lässt sich zeitlich in zwei charakteristische
    Phasen unterteilen.
    Innerhalb der ersten ca. $80\,\text{ms}$ nach dem Direktschall treffen die
    \emph{frühen Reflexionen} (Early Reflections) ein: diskrete, noch einzeln wahrnehmbare
    Echos über klar definierte, kurze Ausbreitungspfade, die das Raumgefühl und die
    Tiefenwahrnehmung maßgeblich prägen.
    Mit zunehmender Zeit überlagern sich Reflexionen höherer Ordnung zu einem dichten,
    statistisch gleichverteilten Nachhallfeld. Diese \emph{späte Nachhallzeit}
    (Late Reverberation) klingt exponentiell ab und wird durch die Nachhallzeit $RT_{60}$
    charakterisiert -- der Zeitraum, in dem der Schalldruckpegel um $60\,\text{dB}$ abfällt.
    Das stochastische Raytracing erfasst beide Phasen natürlich, da sowohl kurze Direktpfade
    als auch komplexe Mehrfachreflexionen in der Impulsantwort akkumuliert werden.


    \subsection{Monte-Carlo-Integration und Raytracing}
    \label{subsec:montecarlo}

    \paragraph{Das akustische Integrationsproblem}

    Der Schalldruck am Empfänger ist das Integral über alle möglichen Schallpfade vom Emitter,
    gewichtet nach Weglänge und materialabhängigen Verlusten. Da die Anzahl dieser Pfade durch
    Mehrfachreflexionen kombinatorisch explodiert, ist eine analytische Auswertung für
    beliebige Raumgeometrien nicht praktikabel. Deterministisches Ray Tracing kann zwar exakte
    Pfade verfolgen, skaliert jedoch exponentiell mit der Anzahl der Reflexionen.

    \paragraph{Der Monte-Carlo-Ansatz}

    Die Monte-Carlo-Integration umgeht dieses Problem durch stochastische Stichprobennahme.
    Statt alle Pfade zu berechnen, schätzt man den Erwartungswert eines Integrals durch den
    arithmetischen Mittelwert von $N$ zufällig gezogenen Stichproben:
    \[\mathbb{E}[f(X)] \approx \frac{1}{N} \sum_{i=1}^{N} f(x_i)\]
    Nach dem Gesetz der großen Zahlen konvergiert dieser Schätzer für $N \to \infty$ gegen
    den wahren Integralwert. Der statistische Fehler nimmt dabei mit $\frac{1}{\sqrt{N}}$ ab:
    ein viermal so großes Stichprobenvolumen halbiert den Fehler.

    \paragraph{Anwendung: Stochastisches Raytracing}

    In der Simulation entspricht eine Stichprobe einem Schallstrahl, der in einer zufälligen
    Richtung $\theta \in [0, 2\pi)$ vom Emitter ausgesendet wird:
    \[\vec{D} = (\cos\theta,\; \sin\theta), \quad \theta \sim \mathcal{U}(0,\; 2\pi)\]
    Jeder Strahl verfolgt seinen Ausbreitungspfad, reflektiert an Wänden und registriert
    einen Treffer, sobald er die Mikrofonkapsel schneidet. Der akkumulierte Beitrag dieses
    Treffers -- Laufzeit $\tau$ und Restdruck $p$ -- repräsentiert eine Stichprobe des
    akustischen Integrals. Über $N = 10^6$ bis $10^8$ Strahlen konvergiert die Gesamtheit
    der Treffer zur vollständigen Impulsantwort des Raumes.
    Der wesentliche Vorteil gegenüber deterministischen Verfahren liegt in der trivialen
    Parallelisierbarkeit: Da jeder Strahl unabhängig von allen anderen ist, lassen sich die
    $N$ Stichproben direkt auf beliebig viele CPU-Kerne aufteilen
    (siehe Kapitel~\ref{sec:concurrency}).


    \subsection{Digitale Signalverarbeitung: Impulsantwort und schnelle Faltung}
    \label{subsec:dsp}

    \paragraph{Die Impulsantwort (IR)}

    Ein lineares, zeitinvariantes System -- wie ein Raum aus akustischer Perspektive -- ist
    vollständig durch seine \emph{Impulsantwort} $h(t)$ charakterisiert: die Ausgabe des
    Systems bei einem idealen Dirac-Impuls $\delta(t)$ am Eingang.
    Die von der Rust-Engine berechnete Menge an (Laufzeit, Druck)-Paaren
    $\{(\tau_i, p_i)\}_{i=1}^{M}$ stellt eine diskrete Approximation dieser Impulsantwort dar:
    \[h(t) \approx \sum_{i=1}^{M} p_i \cdot \delta(t - \tau_i)\]
    Das Python-Modul \texttt{convolution.py} diskretisiert diese Darstellung in ein
    sample-genaues Array, indem jede Verzögerung $\tau_i$ in den Sample-Index
    $n_i = \lfloor \tau_i \cdot f_s \rceil$ umgerechnet und Druckwerte akkumuliert werden
    (bei $f_s = 44100\,\text{Hz}$).

    \paragraph{Das Faltungsintegral}

    Das Ausgangssignal $y(t)$ -- das hallbehaftete Audio -- ergibt sich aus der linearen
    Faltung des Eingangssignals $x(t)$ mit der Impulsantwort $h(t)$:
    \[y(t) = (x * h)(t) = \int_{-\infty}^{\infty} x(\tau)\, h(t - \tau)\, d\tau\]
    In der diskreten Signalverarbeitung wird dieses Integral zur Summe:
    \[y[n] = \sum_{k=0}^{M-1} h[k] \cdot x[n - k]\]
    Die direkte Auswertung hat eine Rechenkomplexität von $\mathcal{O}(N \cdot M)$, wobei $N$
    die Länge des Audiosignals und $M$ die Länge der Impulsantwort bezeichnet.
    Bei realen Impulsantworten mit $M$ von mehreren hunderttausend Samples ist eine direkte
    Berechnung unpraktikabel.

    \paragraph{Schnelle Faltung via FFT (Overlap-Add)}

    Das \emph{Faltungstheorem} der Fouriertransformation besagt, dass eine Faltung im
    Zeitbereich einem punktweisen Produkt im Frequenzbereich entspricht:
    \[Y(f) = \mathcal{F}\{x\}(f) \cdot \mathcal{F}\{h\}(f)
      \implies y[n] = \mathcal{F}^{-1}\{Y(f)\}\]
    Durch den Einsatz der Fast Fourier Transform (FFT) sinkt die Komplexität auf
    $\mathcal{O}((N+M)\log(N+M))$, was bei großen $M$ einen massiven Speedup bedeutet.

    Da das Audiosignal wesentlich länger sein kann als eine einzelne FFT erlaubt,
    wird das \emph{Overlap-Add}-Verfahren eingesetzt: Das Signal $x[n]$ wird in
    überlappungsfreie Blöcke der Länge $L$ aufgeteilt, jeder Block wird separat mit $h[n]$
    gefaltet, und die resultierenden Ausgabeblöcke -- je $L + M - 1$ Samples lang -- werden
    additiv überlagert. Die Implementierung nutzt hierfür
    \texttt{scipy.signal.oaconvolve}, welche diesen Algorithmus effizient realisiert und
    sowohl mono- als auch mehrkanalige Signale verarbeitet.


    \section{Physikalischer Kern: 2D-Strahlensegmentschnitt}
    \label{sec:2d_ray_intersection}

    Der rechnerische Kern der Physik-Engine basiert auf robuster linearer Algebra, um Randkollisionen zu lösen.
    Eine genaue mathematische Modellierung in diesem Stadium ist für hochauflösende akustische Simulationen entscheidend.
    Um den Rechendurchsatz zu maximieren, was bei der Berechnung von Millionen von Strahlen zwingend erforderlich ist, verzichtet die Engine auf
    rechnerisch aufwändige trigonometrische Funktionen zugunsten reiner Vektormathematik.

    \subsection{Geometrische Definitionen}\label{subsec:geometric-definitions}

    %The two primary geometric entities in the 2D simulation space are the acoustic ray and the physical boundary (representing a wall).
    Die beiden wichtigsten geometrischen Entitäten im 2D-Simulationsraum sind der akustische Strahl und die physikalische Grenze (die eine Wand darstellt).

    \paragraph{Der Strahl}
    %A sound ray is defined by an origin point $O$ and a normalized direction vector $\vec{D}$.
    %Any point $P$ along the ray's trajectory can be calculated as a function of the scalar distance $t$ traveled along that direction:
    %\[P(t) = O + t\vec{D}\]
    %where $t \ge 0$ ensures the ray propagates forward in space from the emitter.
    Ein Schallstrahl wird durch einen Ursprungspunkt $O$ und einen normalisierten Richtungsvektor definiert $\vec{D}$.
    Jeder Punkt $P$ entlang der Strahlenbahn kann als Funktion der skalaren Distanz $t$ entlang dieser Richtung berechnet werden:
    \[P(t) = O + t\vec{D}\]
    wobei $t \ge 0$ sicherstellt, dass der Strahl vom Emitter aus im Raum weitergeht.


    \paragraph{Die Wand (Die Strecke)}
    %A wall is represented as a finite line segment defined by two endpoints, $A$ and $B$.
    %Any point $W$ on this segment is found by interpolating between the endpoints using a scalar parameter $u$.
    %Defining the wall vector as $\vec{V} = B - A$, the parametric equation is given by:
    %\[W(u) = A + u\vec{V}\]
    %For a point to lie strictly within the physical bounds of the wall segment, the condition $0 \le u \le 1$ must be satisfied.

    Eine Wand wird als Liniensegment dargestellt, das durch zwei Endpunkte, $A$ und $B$, definiert ist.
    Jeder Punkt $W$ auf diesem Segment wird durch Interpolation zwischen den Endpunkten mit einem skalaren Parameter $u$ gefunden.
    Definiert man den Wandvektor als $\vec{V} = B - A$, ergibt sich die parametrische Gleichung durch:
    \[W(u) = A + u\vec{V}\]
    Damit ein Punkt strikt innerhalb der physikalischen Grenzen des Wandsegments liegt, muss die Bedingung $0 \le u \le 1$ erfüllt sein.

    \subsection{Schnittpunktberechnung}\label{subsec:intersection-formula}

    %To determine the point of intersection, the ray and wall equations are equated:
    %\[O + t\vec{D} = A + u\vec{V}\]
    %This forms a system of linear equations to solve for $t$ (propagation distance) and $u$ (segment intersection parameter).
    %By defining a vector $\vec{E} = A - O$ that points from the ray's origin to the wall's starting point, the equation rearranges to:
    %\[t\vec{D} - u\vec{V} = \vec{E}\]
    %In two-dimensional space, this system can be efficiently solved using the 2D analog of the cross product (often referred to as the perp-dot product).
    %For two arbitrary 2D vectors $\vec{v_1} = (x_1, y_1)$ and $\vec{v_2} = (x_2, y_2)$, this yields a scalar determinant: $\vec{v_1} \times \vec{v_2} = x_1 y_2 - y_1 x_2$.
    %
    %By defining the intersection determinant as $\det = \vec{D} \times \vec{V}$, the parameters $t$ and $u$ can be isolated and solved algebraically:
    %\begin{gather*}
    %    t = \frac{\vec{E} \times \vec{V}}{\det}\\
    %    u = \frac{\vec{E} \times \vec{D}}{\det}\\
    %\end{gather*}

    Um den Schnittpunkt zu bestimmen, werden die Strahlen- und Wandgleichungen gleichgesetzt:
    \[O + t\vec{D} = A + u\vec{V}\]
    Dies bildet ein System linearer Gleichungen, um $t$ (Ausbreitungsdistanz) und $u$ (Segmentschnittparameter) zu lösen.
    Durch die Definition eines Vektors $\vec{E} = A - O$, der vom Ursprung des Strahls zum Ausgangspunkt der Wand zeigt, ordnet sich die Gleichung um:
    \[t\vec{D} - u\vec{V} = \vec{E}\]
    Im zweidimensionalen Raum kann dieses System effizient mit dem 2D-Analogon des Kreuzprodukts gelöst werden (oft als Perp-Dot-Produkt bezeichnet).
    Für zwei beliebige 2D-Vektoren $\vec{v_1} = (x_1, y_1)$ und $\vec{v_2} = (x_2, y_2)$ ergibt sich eine skalare Determinante:$\vec{v_1} \times \vec{v_2} = x_1 y_2 - y_1 x_2$.

    Durch die Definition der Schnittdeterminante als $\det = \vec{D} \times \vec{V}$ können die Parameter $t$ und $u$ isoliert und algebraisch gelöst werden:
    \begin{gather*}
        t = \frac{\vec{E} \times \vec{V}}{\det}\\
        u = \frac{\vec{E} \times \vec{D}}{\det}\\
    \end{gather*}


    \subsection{Schnittpunkt-Validierung (Trefferbedingungen)}\label{subsec:intersection-validation-(hit-conditions)}

    %A mathematically valid solution to the linear system does not inherently guarantee a physical collision.
    %The following constraints must be verified to confirm a valid acoustic reflection:
    %\begin{itemize}
    %    \item \textbf{Parallelism:} If $\det \approx 0$, the ray and the wall are perfectly parallel and will never intersect.
    %    In the implementation, a small epsilon threshold is used rather than an exact zero check to prevent floating-point inaccuracies.
    %    \item \textbf{Directionality:} If $t < 0$, the mathematical lines intersect, but the collision occurs behind the ray's point of emission.
    %    \item \textbf{Segment Bounds:} If $u < 0$ or $u > 1$, the ray intersects the infinite line containing the wall, but misses the physical line segment (e.g., simulating sound passing through an open doorway).
    %\end{itemize}
    %A physical boundary collision is only registered if $\det \neq 0$, $t \ge 0$, and $0 \le u \le 1$.

    Eine mathematisch gültige Lösung des linearen Systems garantiert nicht automatisch eine physikalische Kollision.
    Die folgenden Einschränkungen müssen überprüft werden, um eine gültige akustische Reflexion zu bestätigen:
    \begin{itemize}
        \item \textbf{Parallelität:} Wenn $\det \approx 0$, sind der Strahl und die Wand perfekt parallel und schneiden sich niemals
        In der Implementierung ist ein kleiner Epsilon-Schwellenwert anstelle einer exakten Nullprüfung verwendet, um Gleitkommaungenauigkeiten zu vermeiden:

    \begin{lstlisting}[language=Rust]
    if det.abs() < 1e-6 {
        return None;
    }
    \end{lstlisting}

        \item \textbf{Richtungsabhängigkeit:} Wenn $t < 0$, schneiden sich die mathematischen Geraden, aber die Kollision erfolgt hinter dem Aussendungspunkt des Strahls:
    \begin{lstlisting}[language=Rust]
    if ray_t < 0.0 {
        return None;
    }
    \end{lstlisting}
        \item \textbf{Segmentgrenzen:} Wenn $u < 0$ oder $u > 1$, schneidet der Strahl zwar die unendliche Gerade, auf der die Wand liegt, verfehlt jedoch die physische Strecke
    \begin{lstlisting}[language=Rust]
    if wall_u < 0.0 || wall_u > 1.0 {
        return None;
    }
    \end{lstlisting}
    \end{itemize}

    \subsection{Vektorreflexion}\label{subsec:reflection-vector-calculation}

    %Upon confirming a valid intersection, the resulting trajectory of the sound wave must be calculated.
    %This requires the surface normal vector $\vec{N}$, which points strictly perpendicular away from the wall.
    %For a given 2D wall vector $\vec{V} = (x, y)$, the normalized perpendicular normal is $\vec{N} = (-y, x)$.

    Nach der Bestätigung eines gültigen Schnittpunkts muss die resultierende Trajektorie der Schallwelle berechnet werden.
    Dies erfordert den Oberflächennormalenvektor $\vec{N}$, der streng senkrecht von der Wand weg zeigt
    Für einen gegebenen 2D-Wandvektor $\vec{V} = (x, y)$ ist die normierte senkrechte Normale $\vec{N} = (-y, x)$.

    %The reflected ray direction $\vec{R}$ is determined by taking the incoming ray direction $\vec{D}$ and subtracting twice the projection of that direction onto the surface normal:
    %\[\vec{R} = \vec{D} - 2(\vec{D} \cdot \vec{N})\vec{N}\]
    %Within the Rust codebase, this vector reflection calculation is optimized using the \texttt{reflect} method provided by the \texttt{glam} crate,
    %which leverages SIMD instructions to perform this operation with minimal overhead.

    Die Richtung des reflektierten Strahls $\vec{R}$ wird bestimmt, indem von der Richtung des einfallenden Strahls $\vec{D}$ das Doppelte der Projektion dieser Richtung
    auf die Oberflächennormale subtrahiert wird.
    \[\vec{R} = \vec{D} - 2(\vec{D} \cdot \vec{N})\vec{N}\]
    Innerhalb der Rust-Codebasis wird diese Berechnung der Vektorreflexion durch die \texttt{reflect} Methode optimiert, die vom \texttt{glam}-Crate bereitgestellt wird und
    Single Instruction, Multiple Data (SIMD)-Instruktionen nutzt, um diese Operationen mit minimalem Overhead durchzuführen.



    \section{Hochleistungs-Nebenläufigkeit und Datenparallelität}
    \label{sec:concurrency}

    %A primary advantage of Rust in high-performance computing is its strict compile-time guarantees against data races, colloquially known as
    %``fearless concurrency.'' Distributing a heavily iterative physics loop across multi-core processor architectures in traditional languages often
    %requires complex thread-locking logic (e.g., Mutexes) to prevent memory corruption.
    %This engine avoids locking overhead entirely by employing a functional, data-parallel paradigm utilizing the \texttt{rayon} crate.

    Ein wesentlicher Vorteil von Rust im High-Performance-Computing sind seine strengen Garantien zur Kompilierzeit gegen Data-Races.
    Die Verteilung einer rechenintensiven, iterativen Physikschleife auf Multi-Core-Prozessorarchitekturen erfordert in traditionellen
    Programmiersprachen oft komplexe Thread-Sperrlogiken (z.B. Mutexe), um Speicherbeschädigungen zu verhindern.
    In dieser Engine wird dieser Overhead vermieden durch ein funktionales, datenparalleles Paradigma unter Verwendung des \texttt{rayon}-Crates.

    \subsection{Architektonischer Wandel: Von Shared Mutation zu thread-lokalem-Folding}\label{subsec:architectural-shift:-from-shared-mutation-to-thread-local-folding}

    %In a standard single-threaded execution, sequential iterations append calculated intersection data (e.g., delay and pressure values) into a shared dynamic array.
    %Forcing this loop across numerous threads would result in concurrent write attempts to the same memory addresses, leading to immediate data races.
    %Rather than wrapping these shared arrays in computationally expensive, thread-safe locks, the architecture leverages \texttt{rayon}'s \texttt{into\_par\_iter()} trait
    %combined with a \texttt{fold} and \texttt{reduce} pipeline.

    Bei einer standardmäßigen Single-Thread-Ausführung hängen sequenzielle Iterationen die berechneten Schnittpunktdaten (z.B. delay und pressure werte) an ein gemeinsames dynamisches Array an.
    Diese Schleife auf zahlreiche Threads aufzuteilen, würde zur gleichzeitigen Schreibversuchen auf dieselben Speicheradressen führen, was unmittelbare Data Races zur Folge hätte.
    Anstatt diese gemeinsam genutzt Arrays in rechenintensive, threadsichere Sperren (Locks) zu hüllen, nutzt die Architektur das Trait \texttt{into\_par\_iter()} von \texttt{rayon} in
    Kombination mit einer \texttt{fold}- und \texttt{reduce}- Pipeline.

    Dieser Ansatz teilt die Arbeitslast sicher auf alle verfügbaren CPU-Kerne auf.
    Anstatt zu versuchen, in ein globales Array zu schreiben oder Speicher für jeden Strahl
    dynamisch zu allokieren, verwendet die Engine die \texttt{fold}-Operation, um einen privaten, vorab allokieren, thread-lokalen Puffer für jeden Kern bereitzustellen.

    Jeder Thread simuliert unabhängig seine zugewiesene Teilmenge an Strahlen und hängt die Schnittpunktergebnisse sequenziell in seinen eigenen lokalen Speicherbereich an, ohne Sperren (lock-free).
    Sobald ein Thread seine Arbeitslast abgeschlossen hat, wird die \texttt{reduce}-Phase aktiviert, welche diese isolierten thread-lokalen Puffer sicher und effizient zu den endgültigen,
    vereinheitlichten Arrays zusammenführt.
    Dieser funktionale Ansatz garantiert die Datensicherheit und maximiert gleichzeitig die Hardwareauslastung.

    %This approach safely divides the workload—the generation and simulation of acoustic rays—among all available CPU cores.
    %Instead of attempting to write to a global array or dynamically allocating memory for every individual ray, the engine utilizes the \texttt{fold} operation to provision a private,
    %pre-allocated, thread-local buffer for each core.

    %Each thread independently simulates its assigned subset of rays and sequentially appends the intersection results into its own local memory space, completely lock-free.
    %Once a thread completes its workload, the \texttt{reduce} phase activates, safely and efficiently merging these isolated thread-local buffers into the final, unified arrays.
    %This functional approach guarantees data safety while maximizing hardware utilization.

    \begin{lstlisting}[language=Rust]
use rayon::prelude::*;

fn run_simulation_parallel(
config: &SimulationConfig,
walls: &Vec<Wall>
) -> (Vec<f32>, Vec<f32>) {
    let(delays_par, pressures_par, _) = (0..config.rays_to_cast)
        .into_par_iter()
        // 1. Falten, Split
        .fold( || {
            let local_delays = Vec::with_capacity(CAPACITY);
            let local_pressures = Vec::with_capacity(CAPACITY);
            let temp_ray_buffer = Vec::with_capacity(CAP);

            (local_delays,local_pressures,temp_ray_buffer)
        },
            |mut local_buf, _| {
                let fresh_ray = generate_initial_ray(config);
                simulate_single_ray(fresh_ray,walls,config,&mut local_buf.2);
                for(delay, pressure) in &local_buf.2 {
                    local_buf.0.push(*delay);
                    local_buf.1.push(*pressure);
                }
                local_buf
            }
        // 2. Reduzieren, Global Merge
        ).reduce(
        || (Vec::new(),Vec::new(),Vec::new()),
        |mut a, mut b| {
            a.0.append(&mut b.0);
            a.1.append(&mut b.1);
            a
        }
    );
    (delays_par, pressures_par)
}
    \end{lstlisting}


    \section{Leistungsoptimierungen und Profiling}
    \label{sec:optimizations}

    Benchmarking the compiled binary (utilizing Rust's \texttt{--release} flag for maximum LLVM optimization)
    revealed significant bottlenecks and subsequent optimization opportunities within the parallelized architecture.

    Das Benchmarking der kompilierten Binärdatei (unter Verwendung des \texttt{--release} flag und dem \texttt{let var = Instant::now();} timer) deckte signifikante
    Engpässe und darauffolgende Optimierungsmöglichkeiten innerhalb der parallelisierten Architektur auf.


    \subsection{Reduzierung des Speicherallokations-Overheads: \texttt{flat\_map} vs. \texttt{fold}}
    \label{sec:memory_overhead}


    Während die anfängliche parallele Implementierung den \texttt{flat\_map}-iterator von \texttt{rayon} nutzt, offenbarte die Skalierung der Simulation auf $100.000.000$ Strahlen einen
    signifikanten Leistungsengpass bezüglich der Speicherallokation.
    Die \texttt{flat\_map}-Operation allozierte dynamisch einen neuen \texttt{Vec} für die Rückgabedaten jedes einzelnen Strahls.
    Bei der Verteilung auf alle CPU-Kerne führte dies zu Millionen von flüchtigen Heap-Allokationen pro Sekunde.
    Das System erlebte Memory Thrashing; die CPU-Auslastung fiel auf 0\% der Hauptspeicher war ausgelastet, und das Betriebssystem entzog dem Ausführungs-Thread die Ressourcen,
    was zum Stillstand der Simulation führte.

    Um diesen Allokations-Overhead zu beheben, wurde die parallele Pipeline refaktorisiert, um das \texttt{fold}-Muster zu verwenden.
    Anstatt Speicher pro Strahl zu allokieren, stellt die \texttt{fold}-Operation exakt einen thread-lokalen Puffer (initialisiert über \texttt{Vec::with\_capacity();}) pro CPU-Kern bereit.
    Die Kerne akkumulieren ihre lokalen Ergebnisse und führen diese erst nach Abschluss des Threads zusammen.
    Diese architektonische Änderung reduzierte die Heap-Allokationen um mehrere Größenordnungen und ermöglichte es dem System, $100.000.000$ Strahlen effizient zu verarbeiten,
    ohne die Speicherressourcen zu erschöpfen.
    %While the initial parallel implementation utilized \texttt{rayon}'s \texttt{flat\_map} iterator, scaling the simulation to $100,000,000$ rays exposed a critical flaw in memory allocation.
    %The \texttt{flat\_map} operation dynamically allocated a new \texttt{Vec} for every single ray's return data.
    %When distributed across all CPU cores, this resulted in millions of transient heap allocations per second.
    %The system experienced severe memory thrashing; CPU utilization dropped to 0\%, system memory was exhausted, and the OS starved the execution thread, causing the simulation to stall.
    %
    %To resolve this allocator crisis, the parallel pipeline was refactored to use the \texttt{fold} pattern.
    %Instead of allocating memory per ray, the \texttt{fold} operation provisions exactly one thread-local buffer (initialized via \texttt{Vec::with\_capacity()}) per CPU core.
    %The cores accumulate their local results and merge them only upon thread completion.
    %This architectural shift reduced heap allocations by orders of magnitude, allowing the system to process $100,000,000$ rays without memory exhaustion.


    \subsection{Benchmarking Ergebnisse und I/O Overhead}
    \label{sec:benchmarks}

    Empirische Tests verdeutlichen die Wirksamkeit der optimierten Parallelisierung sowie die verborgenen Kosten von System-I/O und Thread-Initialisierung.
    %Empirical testing highlighted the effectiveness of the optimized parallelization, as well as the hidden costs of system I/O and thread initialization.

    \begin{itemize}
        \item \textbf{Paralleler Speedup:} Bei einer moderaten Arbeitslast wurde die Single-Thread-Simulation in $11,79$ Sekunden abgeschlossen.
        Die durch \texttt{rayon} optimierte parallele Simulation bewältigt dieselbe Arbeitslast in $2,56$ Sekunden, was einen $4{,}60\times$ Speedup auf der Test-Hardware ergab.
        \item \textbf{Thread-Pool-Kaltstart:} Bei extrem kleinen Arbeitslasten wies \texttt{rayon} einen Initialisierung-Overhead auf.
        In einem Benchmark mit wenigen Strahlen wurde die Single-Thread-Ausführung in $0,14$ Sekunden abgeschlossen, während die parallele Ausführung $0,34$ Sekunden in Anspruch nahm.
        Dies deutet darauf hin, dass die Kosten für die Initialisierung des Thread-Pools die Recheneinsparungen bei trivialen Arbeitslasten überwiegen
        \item \textbf{The I/O Engpass:} Das Profiling der End-to-End-Ausführungszeit deckte einen massiven Engpass bei der Datenserialisierung auf.
        In einem Test mit $10.000.000$ Strahlen wurde die interne Simulation schnell ausgeführt($4,72s$ Single-Thread vs. $0,64s$ Parallel).
        Die anschließende Serialisierung und das Schreiben des Datensatzes dominierten jedoch die End-to-End-Laufzeit($37,40s$ vs. $33,24s$),
        was die Notwendigkeit eines RAM-gestützten \texttt{BufWriter} unterstreicht.
    \end{itemize}
    %\begin{itemize}
    %    \item \textbf{Parallel Speedup:} For a moderate workload, the single-threaded simulation completed in 11.79 seconds. The \texttt{rayon}-optimized parallel simulation completed the same workload in 2.56 seconds, yielding a $4.60\times$ speedup on the test hardware.
    %    \item \textbf{Thread Pool Cold Start:} For extremely small workloads, \texttt{rayon} exhibited an initialization overhead.
    %    In a low-ray benchmark, the single-threaded execution completed in 0.14 seconds, whereas the parallel execution took 0.34 seconds.
    %    This indicates that the cost of initializing the thread pool outweighs the computational savings for trivial workloads.
    %    \item \textbf{The I/O Bottleneck:} Profiling the end-to-end execution time revealed a massive bottleneck in data serialization.
    %    In a test generating 1,533,997 acoustic hits, the internal simulation executed rapidly, but writing the massive dataset to \texttt{ir\_output.json} took over 30 seconds (32.27s single-threaded total vs.
    %    31.28s parallel total).
    %    This demonstrated that standard file I/O operations entirely eclipsed the physics engine's execution time,
    %    necessitating the subsequent implementation of a \texttt{BufWriter} to stage the serialization in RAM.
    %\end{itemize}
\end{document}
